\documentclass[10pt]{article}
\usepackage[left=0.75in,top=0.75in,right=0.75in,bottom=0.75in]{geometry}
\usepackage{enumitem}
\usepackage{hyperref}
\usepackage{xcolor}
\usepackage{titlesec}

\definecolor{darkblue}{RGB}{0,51,102}
\hypersetup{colorlinks=true,linkcolor=darkblue,filecolor=darkblue,urlcolor=darkblue}
\titleformat{\section}{\large\bfseries\color{darkblue}}{}{0em}{}[\titlerule]
\titlespacing{\section}{0pt}{12pt}{6pt}
\setitemize[itemsep=2pt,parsep=0pt]

\begin{document}
\begin{center}
{\LARGE\bfseries Luis Eduardo González González}\\
Ciudad de México | 55 6707 6720 | fectda@gmail.com\\
LinkedIn | GitLab | GitHub
\end{center}

\section{RESUMEN PROFESIONAL}
Profesional de TI con experiencia en full-stack, arquitectura en nube y gestión de proyectos. Especializado en Python, JavaScript y GCP. Enfocado en proyectos independientes: orquestación de IA local, observabilidad y sistemas distribuidos. Base analítica sólida en Ciencias Físicas.

\section{HABILIDADES}
\begin{itemize}
    \item Lenguajes: Python, JavaScript, Java, Go, Fortran, C, PHP
    \item Front-End: HTML5, CSS3, Vue.js, jQuery, Bootstrap
    \item Back-End: Python (Flask), PHP (Laravel), Node.js
    \item IA/Automation: Prompt Engineering, AI Agents, Ollama, n8n
    \item DevOps/Cloud: GCP, AWS, Docker, K8s, Terraform, Ansible
    \item Observabilidad: Grafana, Istio, Prometheus
    \item Bases de Datos: PostgreSQL, MySQL, MongoDB, Redis
\end{itemize}

\section{EXPERIENCIA}
\subsection{Solutions Architect \& Maker | Independiente}
\textit{2026 - Actualidad}
\begin{itemize}
    \item Infraestructura on-premise: Xen Orchestra, K8s, RabbitMQ
    \item LLMs locales (Ollama) para casos educativos/creativos
    \item Maker: Hardware personalizado (3D, IoT)
    \item Observabilidad: Grafana, Istio, Prometheus
    \item Seguridad: Cloudflare Tunnel, Zero Trust
\end{itemize}

\subsection{Responsable R\&D | GET IN}
\textit{2024 - 2026}
\begin{itemize}
    \item Pruebas de concepto para viabilidad técnica
    \item Modelo estadístico: 90\% reducción en revisión de datos
\end{itemize}

\subsection{CIO | GET IN}
\textit{2018 - 2024}
\begin{itemize}
    \item Liderazgo de 5 equipos multidisciplinaryes
    \item Migración a GCP (PaaS): reducción de costos, mayor fiabilidad
    \item Reingeniería a microservicios
    \item Procesamiento de datos: 3x velocidad, 30\% reducción costos
    \item Tecnologías: Go, Python, PHP, MongoDB, Vue.js, GCP
\end{itemize}

\subsection{Developer Senior | elcultivo.mx}
\textit{2015 - 2018}
\begin{itemize}
    \item 20 aplicaciones web para 20 clientes en un año
    \item Arquitectura MVC con Laravel y Vue.js
\end{itemize}

\section{PROYECTOS}
Infraestructura (K8s, Terraform, Ansible), IA Generativa (Ollama, n8n), Simulaciones (Basic, Fortran), Redes Neuronales (TensorFlow), Domótica (Home Assistant), 3D (SketchUp, FDM)

\section{EDUCACIÓN}
Maestría en Física | UNAM | 2008-2010\\
Licenciatura en Física | UNAM | 2003-2007

\section{IDIOMAS}
Español: Nativo | Inglés: Avanzado (C1)

\end{document}